\documentclass[11p,a4paper]{article}
\usepackage{graphics}
\usepackage{epsfig}
\usepackage{amsfonts}
\usepackage{latexsym}
\usepackage{psfig}
\newcommand{\N}{{\mathbb N}}
\newcommand{\Z}{{\mathbb Z}}
\newcommand{\Q}{{\mathbb Q}}
\newcommand{\R}{{\mathbb R}}
\newcommand{\C}{{\mathbb C}}

\begin{document}

\noindent
{\Large \bf Symbolic Computation of Formal  Solutions
for 2 and 3 Dimensional Dynamical Systems}
\\

\noindent
G\'erard Eichenm\"uller\footnote{The work of the author has 
been supported by the Gottlieb Daimler- und Karl Benz-Stiftung}, \\
\noindent
LMC--IMAG, B. P. 53, 51 Rue des 
Math\'{e}matiques, 38041 Grenoble Cedex 9, France, 
{\tt gerard.eichenmueller@imag.fr}
\\

\begin{center}
{\large \bf Abstract}
\end{center}

\noindent
First we will consider 2-dimensional autonomous dynamical systems 
\begin{equation}\label{sys}
	\dot X = F(X),\ X={x \choose y},\ F={f \choose g},\ f,g\in \C[x,y]
\end{equation}
and present the implementation of an algorithm proposed by A. Bruno 
that can be used
to find parametrized formal solutions $(x(t),y(t))$ 
in the neighbourhood of the origin 
for these systems if the origin is a simple point or an isolated singularity.
Then we will give some examples how this algorithm can be extended
for the integration of 3-dimensional dynamical systems.
\\

The algorithm is split into three parts according to the classification 
of the dynamical systems into 3 cases:
\begin{itemize}
\item    the origin is a simple point if $F(0)\neq 0$. 
\item    it is an elementary singular point if $F(0)= 0$ and the linear 
part is not nilpotent.
\item    it is a nonelementary singular point if $F(0)= 0$ and the linear 
part is nilpotent.
\end{itemize}
In the case of a simple point we can apply the flow box theorem to transform 
the initial system into 
\[
     \dot {\tilde X}= {1 \choose 0}
\]
and use Taylor series to find the transformation $X=H(\tilde X)$.

If the origin is an elementary singular point we can apply normal 
form constructions to
reduce the initial system to an integrable system. We compute the 
Poincar\'e-Dulac normal form if the eigenvalues $\lambda_1$ and 
$\lambda_2$ of the linear part are in $\R$ or if we have complex 
coefficients in F. If we have only rational coefficients in F and 
$\lambda_1,\lambda_2 \in \C$ we calculate the rational normal form 
which is linear, 
if $\lambda_1$ and $\lambda_2$ are purely imaginary we use 
the rational normal form and polar coordinates.

In the case of a nonelementary singular point the Newton polygon is used to
divide the neighbourhood of the origin into sectors. Now we calculate formal
solutions valid on these sectors. For all sectors corresponding to 
vertices of the Newton polygon we apply a new time transformation and a 
power transformation to transform the initial system to an system that 
has an elementary singular point in the origin. On the sectors 
corresponding to edges of the Newton polygon we use 
directional blowing up to desingularize the singularity. This  
leads to a finite number of systems that can be integrated using the 
methods described above. 

The program has been implemented in Maple. It can be used and tested online at
{\tt http://www-lmc.imag.fr/lmc-cf/Gerard.Eichenmueller/Software/}. 
We will treat some examples to illustrate each of the possible cases and 
to show how the algorithm works.

This algorithm can be applied to integrate 3-dimensional dynamical 
system 
\[
\dot X = F(X),\ X=\left( \begin{array}{c}x\\y\\z \end{array} \right),\ 
F=\left( \begin{array}{c}f\\g\\h \end{array} \right),\ f,g,h \in \C[x,y,z]
\]
if the origin is an elementary singular point and the 
Poincar\'e-Dulac normal form 
\[
\left\{ \begin{array}{lrl}\frac{dx}{dt}= & \lambda_1 x + \sigma_1 y & 
+ x \sum_{\langle Q,\Lambda \rangle=0} a_Q X^Q\\
\frac{dy}{dt}= & \lambda_2 y + \sigma_2 z & + y 
\sum_{\langle Q,\Lambda \rangle=0} b_Q X^Q \\
\frac{dz}{dt}= & \lambda_3 z & + z \sum_{\langle Q,\Lambda \rangle=0} 
c_Q X^Q 
\end{array}\right.
\]
has its support on a plane. That means that the resonance equation 
\[
\langle Q, \Lambda \rangle = 0
\]
has 2 linearly independent solutions in $\N^3$. Then the 
3-dimensional system can be reduced to a 2-dimensional system using 
a change of coordinates 
described by a power transformation. The resulting sytem has a nonlementary 
singularity in 
the origin and we can find a division of the neighbourhood into sectors 
and formal solution on those sectors. The two-dimensional solutions 
can be retransformed into $(x,y,z)$ coordinates. But the approbiate 
power transformation that reduces the system does not always exist. 
We will show when it is possible to find it
and how to construct it if it exists.

\begin{thebibliography}{99}

\bibitem{bruno1}
	A. Bruno {\sf Local methods in non linear differential equations}
	Springer Verlag 1989.

\bibitem{jean}
	J. Della Dora {\sf Introduction to normal form theory}, Preprint.

\bibitem{stolo}
	L. Stolovitch {\sf Syst\`emes dynamiques et formes normales, rapport 
	de recherche}, technical report, INPG, 1992.

\bibitem{vallier}
	L. Vallier {\sf An Algorithm for the computation of Normal Forms and 
	invariant Manifolds}, ISAAC, page 225-233, edit. Manuel Bronstein, 
	1993.

\bibitem{pelletier}
M.~Pelletier.
{\sf Eclatements quasi-homog\`enes.}
Technical Report 92-04, Universit\'e de Bourgogne, 1992.

\bibitem{hale}
Jack K.~Hale, Shui-Nee~Chow.
{\sf Methods of Bifurcation Theory}.
Springer-Verlag, 1986.

\end{thebibliography}


\end{document}